\documentclass[11pt]{article}
\usepackage[T1]{fontenc}
\usepackage[utf8]{inputenc}
\usepackage{lmodern,amsmath,amssymb,amsthm,mathtools,booktabs}
\usepackage[margin=1in]{geometry}
\usepackage[colorlinks=true,linkcolor=blue,urlcolor=blue,citecolor=blue]{hyperref}
\usepackage{microtype}
\newtheorem{theorem}{Theorem}
\newtheorem{lemma}{Lemma}
\newtheorem{proposition}{Proposition}
\newcommand{\coef}[2]{[#1^{#2}]}
\newcommand{\fall}[2]{(#1)_{#2}}
\newcommand{\Z}{\mathbb Z}
\newcommand{\N}{\mathbb N}
\DeclareMathOperator{\supp}{supp}
\DeclareMathOperator{\e}{e}
\title{Sparse nonattacking chess placements\\\large Every fixed asymptotic order and growth-index inversion}
\author{}
\date{2 October 2026}
\begin{document}
\maketitle
\begin{abstract}
The numbers of ways to place $n$ nonattacking kings or knights on an $n\times n$ board have the same classical leading equivalent, $e^{-9/2}n^{2n}/n!$. Their first relative corrections are different: $7/(3n)$ for kings and $37/(3n)$ for knights. We compute four corrections exactly and give a finite algorithm and a uniform remainder proof for every fixed order. The method applies to every fixed finite symmetric move set, including a uniform extension with $k/n$ in a compact positive interval. Exact boundary contributions and triangle density explain the first correction. We also derive a Lambert-$W$ growth-index inverse and state the necessary integer-rounding qualification. The leading law and general cluster-expansion methods are prior work; the results presented here are explicit specialized expansions with controlled errors.
\end{abstract}

\section{Definitions and main results}
Let $K_n$ count subsets of $n$ squares of the ordinary $n\times n$ board containing no attacking pair of kings, and let $N_n$ be the corresponding knight count. These are OEIS \href{https://oeis.org/A201513}{A201513} and \href{https://oeis.org/A201540}{A201540}. Rotations and reflections remain distinct. There is no one-piece-per-row or one-piece-per-column constraint.

Kot\v e\v sovec's leading equivalent, dated 29 November 2011 and reproduced in his book~\cite{Kote}, is
\begin{equation}\label{eq:classical}
 K_n\sim N_n\sim B_n,\qquad B_n=e^{-9/2}\frac{n^{2n}}{n!}.
\end{equation}
It is not a new result of this report. The book also contains fixed-piece polynomial formulas. Pages 77 and 293 display the diagonal leading laws; pages 685--686 give the general bounded-move leading argument. A fixed-piece polynomial expansion alone does not justify substituting a growing piece count in its remainder. We supply that uniform step below.

\begin{theorem}\label{thm:main}
Both sequences have rational relative expansions to every fixed order, with the corresponding next-power remainder. In particular,
\begin{align}
\frac{K_n}{B_n}
 &=1+\frac7{3n}+\frac{119}{9n^2}+\frac{11327}{810n^3}
       +\frac{90475}{1944n^4}+O(n^{-5}),\label{eq:kings}\\
\frac{N_n}{B_n}
 &=1+\frac{37}{3n}+\frac{914}{9n^2}+\frac{330917}{810n^3}
       +\frac{966583}{1944n^4}+O(n^{-5}).\label{eq:knights}
\end{align}
Consequently,
\begin{equation}\label{eq:ratio}
\frac{N_n}{K_n}=1+\frac{10}{n}+\frac{65}{n^2}
 +\frac{332}{3n^3}-\frac{4841}{6n^4}+O(n^{-5}).
\end{equation}
\end{theorem}
The logarithmic corrections are particularly simple:
\begin{align}
\log\frac{K_n}{B_n}
 &=\frac7{3n}+\frac{21}{2n^2}-\frac{379}{30n^3}
 -\frac{357}{40n^4}+O(n^{-5}),\label{eq:klog}\\
\log\frac{N_n}{B_n}
 &=\frac{37}{3n}+\frac{51}{2n^2}-\frac{6559}{30n^3}
 -\frac{1397}{40n^4}+O(n^{-5}).\label{eq:nlog}
\end{align}
If the carrier is instead $e^{-9/2}(en)^n/\sqrt{2\pi n}$, the first relative corrections are $9/(4n)$ and $49/(4n)$. The carrier must therefore be specified when quoting a coefficient.

The mechanism is an application of established cluster-expansion ideas. Canonical independent-set cluster expansions and convergent tails are developed, for example, by Davies, Jenssen, and Perkins~\cite{DJP}, who credit earlier canonical work of Pulvirenti and Tsagkarogiannis~\cite{PT}. We do not claim a new general cluster method, or exhaustive historical priority for each specialized coefficient. The proof here is self-contained and the coefficient computation is exact.

\section{The finite-range theorem and its geometry}
Fix a nonempty finite symmetric set $S\subset\Z^2\setminus\{0\}$, and write $d=|S|$. Let $G_n$ have vertex set $V_n=\{0,\ldots,n-1\}^2$, with $x,y$ adjacent when $y-x\in S$. Write
\begin{equation}\label{eq:logI}
 I_n(z)=\sum_{k\ge0}a(n,k)z^k,\qquad
 \log I_n(z)=\sum_{j\ge1}b_j(n)z^j.
\end{equation}
Here $a(n,k)$ counts independent $k$-subsets. We will prove that for each fixed $j$ there are rational numbers $\alpha_j,\beta_j,\gamma_j$ for which
\begin{equation}\label{eq:quadratic}
 b_j(n)=\alpha_j n^2+\beta_j n+\gamma_j
\end{equation}
for all sufficiently large integers $n$. In particular
\[
(\alpha_1,\beta_1,\gamma_1)=(1,0,0),\qquad
\alpha_2=-\frac{d+1}{2}.
\]
Put $a=\alpha_2$.

\begin{theorem}\label{thm:general}
For every fixed $M\ge0$,
\begin{equation}\label{eq:diagonal-general}
 a(n,n)=\frac{n^{2n}}{n!}e^a
 \left(\sum_{r=0}^M\frac{c_r}{n^r}+O_{M,S}(n^{-M-1})\right),
\end{equation}
where $c_0=1$ and every $c_r$ is rational. A finite formula for $c_r$ uses only the cluster polynomials $b_2,\ldots,b_{r+2}$.

More generally, for any fixed compact interval $J\subset(0,\infty)$, uniformly over integers $k$ with $\theta=k/n\in J$,
\begin{equation}\label{eq:density-general}
 a(n,k)=\frac{n^{2k}}{k!}e^{a\theta^2}
 \left(\sum_{r=0}^M\frac{c_r(\theta)}{n^r}
       +O_{M,S,J}(n^{-M-1})\right).
\end{equation}
The coefficient functions have the explicit finite formula in Section~\ref{sec:transfer}.
\end{theorem}
The actual value $\theta=k/n$ is used in~\eqref{eq:density-general}. Thus a choice $k=\lfloor\lambda n\rfloor$ retains its fractional-part corrections. No uniformity as $\theta\to0$ or $\theta\to\infty$ is asserted.

\subsection{Boundary loss and triangles}
Define
\[
 B=\frac12\sum_{u\in S}(|u_1|+|u_2|),\qquad
 \tau=\frac16\#\{(u,v)\in S^2:v-u\in S\}.
\]
Thus $\tau$ is the number of triangles per bulk vertex, when each triangle is counted once. The first correction is
\begin{equation}\label{eq:first-geometric}
\boxed{c_1(1)=B+\frac13-\tau.}
\end{equation}
For kings $(d,B,\tau)=(8,6,4)$; for knights $(d,B,\tau)=(8,12,0)$. The degree determines the common leading term, but the boundary loss and triangles distinguish the next term.

To verify~\eqref{eq:first-geometric}, the elementary graph identities are
\begin{align*}
 b_2&=-\frac{|V|}{2}-|E|,\\
 b_3&=\frac{|V|}{3}+2|E|+\sum_{v\in V}\binom{\deg(v)}2
       -\#\{\text{triangles}\}.
\end{align*}
Since
\[
 |E(G_n)|=\frac12\sum_{u\in S}(n-|u_1|)(n-|u_2|)
\]
for sufficiently large $n$, these imply
\begin{equation}\label{eq:first-clusters}
 \beta_2=B,\qquad
 \alpha_3=\frac13+\frac{d(d+1)}2-\tau.
\end{equation}
The coefficient extraction in Section~\ref{sec:transfer} then gives the stronger identity
\begin{equation}\label{eq:c1theta}
 c_1(\theta)=\frac{d+1}{2}\theta+B\theta^2
 -\left(\frac{d+1}{2}+\tau-\frac13\right)\theta^3.
\end{equation}

\section{Exact connected-support polynomials}
For a finite graph $H$, write $L_j(H)=[z^j]\log I_H(z)$. For a nonempty set of vertices $T$ define
\begin{equation}\label{eq:weight}
 w_j(T)=\sum_{U\subseteq T}(-1)^{|T|-|U|}L_j(G_n[U]).
\end{equation}
The empty graph contributes zero. Then
\begin{equation}\label{eq:connected}
 w_j(T)=0\quad\text{if }|T|>j\text{ or }G_n[T]\text{ is disconnected},
 \qquad b_j(n)=\sum_{\varnothing\ne T\subseteq V_n}w_j(T).
\end{equation}
Indeed, give each vertex its own indeterminate in the independence polynomial. A degree-$j$ monomial of the formal logarithm uses at most $j$ variables. M\"obius inversion~\eqref{eq:weight} selects precisely the monomials whose support equals $T$, after all variables are equated. If the induced graph on $T$ is disconnected, its independence polynomial factors over components, so its logarithm contains no monomial using all of $T$. This proves~\eqref{eq:connected}.

Enumerate nonempty finite connected lattice sets $T$ of size at most $j$ modulo translations. Normalize each representative so that both coordinate minima are zero; do not quotient by rotations or reflections. Let
\[
 w(T)=\max_{(x,y)\in T}x,\qquad h(T)=\max_{(x,y)\in T}y.
\]
Its number of translates fitting in $V_n$ is exactly $(n-w(T))_+(n-h(T))_+$. Hence the following identity holds for every $n$:
\begin{equation}\label{eq:clipped}
 b_j(n)=\sum_T w_j(T)(n-w(T))_+(n-h(T))_+.
\end{equation}
There are finitely many supports because the move set is fixed. If
\[
 R=\max_{u\in S}\max(|u_1|,|u_2|),
\]
connectedness implies $w(T),h(T)\le R(j-1)$. For $n\ge R(j-1)+1$, equation~\eqref{eq:clipped} is the exact quadratic~\eqref{eq:quadratic}, with
\begin{equation}\label{eq:alphabetagamma}
 \alpha_j=\sum_Tw_j(T),\quad
 \beta_j=-\sum_Tw_j(T)(w(T)+h(T)),\quad
 \gamma_j=\sum_Tw_j(T)w(T)h(T).
\end{equation}
These are finite rational sums, not numerical fits or polynomial interpolation.

\section{Uniform analytic control}
A uniform bound is needed to turn the finite-order data into an expansion for a growing number of pieces.

\begin{lemma}\label{lem:zerofree}
For every graph $G$ of maximum degree at most $d$, its independence polynomial has no zero on
\[
 |z|\le\rho:=\frac1{4(d+1)}.
\]
The logarithm normalized by $\log I_G(0)=0$ satisfies
\begin{equation}\label{eq:bound}
 |\log I_G(z)|\le4|V(G)||z|,\qquad
 |[z^j]\log I_G(z)|\le4|V(G)|\rho^{1-j}.
\end{equation}
\end{lemma}
\begin{proof}
Put $\delta=1/[2(d+1)]$. Induct on the number of vertices, proving that $I_G/I_{G-v}$ exists and differs from $1$ in modulus by at most $\delta$. The deletion identity gives
\[
 \frac{I_G}{I_{G-v}}=1+z\frac{I_{G-N[v]}}{I_{G-v}}.
\]
Delete the at most $d$ neighbors of $v$ one at a time in $G-v$. The induction hypothesis bounds the second quotient by $(1-\delta)^{-d}<2$, since
\[
 d[-\log(1-\delta)]\le\frac{d\delta}{1-\delta}
 =\frac{d}{2d+1}<\frac12.
\]
Thus the ratio differs from $1$ by at most $2|z|\le\delta<1$, closing the induction and proving nonvanishing. Its analytic logarithm has modulus at most
\[
 \frac{2|z|}{1-\delta}\le4|z|.
\]
Telescope over a vertex deletion order. Cauchy's coefficient estimate on $|z|=\rho$ proves the second bound.
\end{proof}

Fix $R_0>1$ and a compact positive interval for $\theta$. The omitted tail after cluster order $J$ in $\log I_n(\theta t/n)$, uniformly on $|t|\le R_0$, is bounded by
\begin{equation}\label{eq:logtail}
 \sum_{j>J}4n^2\rho^{1-j}(\theta R_0/n)^j=O(n^{1-J}).
\end{equation}
Use~\eqref{eq:quadratic} for the finitely many retained indices. Taking $J=M+2$ and exponentiating gives
\begin{equation}\label{eq:Hexp}
 I_n(\theta t/n)e^{-\theta nt}
 =e^{a\theta^2t^2}
 \left(\sum_{r=0}^M H_r(t;\theta)n^{-r}+O(n^{-M-1})\right)
\end{equation}
uniformly on the same disk. Exponentiation preserves the order because the retained expressions are uniformly bounded there. Section~\ref{sec:transfer} identifies $H_r$ and proves that coefficient extraction preserves this error.

\section{A finite all-order extraction formula}\label{sec:transfer}
For $\ell\ge1$ set
\begin{equation}\label{eq:h}
 h_\ell(t;\theta)=\alpha_{\ell+2}(\theta t)^{\ell+2}
   +\beta_{\ell+1}(\theta t)^{\ell+1}
   +\gamma_\ell(\theta t)^\ell,
\end{equation}
with missing indices below $2$ interpreted as zero. Define polynomials $H_r$ by
\[
 \sum_{r\ge0}H_r(t;\theta)\epsilon^r
 =\exp\left(\sum_{\ell\ge1}h_\ell(t;\theta)\epsilon^\ell\right).
\]
Equivalently, $H_0=1$ and
\begin{equation}\label{eq:Hrec}
 rH_r=\sum_{\ell=1}^r\ell h_\ell H_{r-\ell}.
\end{equation}
Define rational polynomials $P_q(m)$ by
\begin{equation}\label{eq:P}
 \sum_{q\ge0}P_q(m)\epsilon^q
 =\exp\left(-\sum_{\ell\ge1}\frac{\epsilon^\ell}{\ell}
                      \sum_{i=0}^{m-1}i^\ell\right).
\end{equation}
Faulhaber's formula gives $\deg P_q\le2q$. The first three are
\[
 P_0=1,\quad P_1=-\frac{m(m-1)}2,\quad
 P_2=\frac{m(m-1)(m-2)(3m-1)}{24}.
\]
Write $D=t\,d/dt$.

\begin{proposition}\label{prop:algorithm}
The coefficients in Theorem~\ref{thm:general} are
\begin{equation}\label{eq:cformula}
 c_r(\theta)=e^{-a\theta^2}
 \sum_{q=0}^r\theta^{-q}
 \left. P_q(D)\{e^{a\theta^2t^2}H_{r-q}(t;\theta)\}\right|_{t=1}.
\end{equation}
In particular $c_r=c_r(1)\in\mathbb Q$.
\end{proposition}
\begin{proof}
For $k=\theta n$ let $H_{n,k}(t)=I_n(\theta t/n)e^{-kt}$. Expanding $e^{kt}$ and extracting $[t^k]$ gives the exact identity
\begin{equation}\label{eq:exacttransfer}
 \frac{a(n,k)}{n^{2k}/k!}
 =\sum_{m=0}^k[t^m]H_{n,k}(t)\frac{\fall{k}{m}}{k^m}.
\end{equation}
The scale in $I_n(\theta t/n)$ is $(k/n^2)^k$, which accounts for the normalization.

For an analytic function $f$ on $|t|\le R_0$, Cauchy's estimate and $0\le\fall{k}{m}/k^m\le1$ show that the functional on the right of~\eqref{eq:exacttransfer} has norm at most $(1-1/R_0)^{-1}$ in the disk supremum norm. Thus the analytic remainder in~\eqref{eq:Hexp} remains $O(n^{-M-1})$ after extraction.

It remains to expand the finite falling product. For $0\le m\le\sqrt{k}$, the elementary symmetric functions in $0,1,\ldots,m-1$ give
\[
 \frac{\fall{k}{m}}{k^m}
 =\sum_{q=0}^M P_q(m)k^{-q}+E_{M,m,k}.
\]
Put $A=m(m-1)/2$. Since the degree-$q$ elementary symmetric function is at most $A^q/q!$,
\begin{equation}\label{eq:falltail}
 |E_{M,m,k}|
 \le \frac{k^{-M-1}A^{M+1}}{(M+1)!}e^{A/k}
 \le C_M k^{-M-1}m^{2M+2}.
\end{equation}
The coefficients of each analytic function in~\eqref{eq:Hexp} decay geometrically, so summing~\eqref{eq:falltail} gives $O(k^{-M-1})$. For $m>\sqrt{k}$ the exact factor is bounded by $1$ up to $k$ and vanishes beyond $k$, while $P_q(m)$ has fixed polynomial growth. These tails, including the extension of each polynomial sum to all $m$, are $O(\operatorname{poly}(k)R_0^{-\sqrt{k}})$. Finally,
\[
 \sum_{m\ge0}m^j[t^m]f(t)=D^jf(1).
\]
Apply this identity to the retained terms and use $k^{-q}=\theta^{-q}n^{-q}$. Since $k$ and $n$ are uniformly comparable,~\eqref{eq:cformula} and the uniform remainder follow. Exponential factors cancel in~\eqref{eq:cformula}; the coefficients are rational at $\theta=1$.
\end{proof}
Theorem~\ref{thm:general}, and hence the existence assertion of Theorem~\ref{thm:main}, is now proved. Inserting~\eqref{eq:first-clusters} into the first operator gives~\eqref{eq:c1theta} and~\eqref{eq:first-geometric}.

\section{Exact coefficients and reproducibility}
For kings and knights the quadratic cluster coefficients through $j=6$ are as follows. Each row means $b_j(n)=\alpha_jn^2+\beta_jn+\gamma_j$ once $n\ge R(j-1)+1$.
\begin{center}
\begin{tabular}{crrr|rrr}
\toprule
&\multicolumn{3}{c}{Kings}&\multicolumn{3}{c}{Knights}\\
$j$&$\alpha_j$&$\beta_j$&$\gamma_j$&$\alpha_j$&$\beta_j$&$\gamma_j$\\
\midrule
2&$-9/2$&6&$-2$&$-9/2$&12&$-8$\\
3&$97/3$&$-76$&44&$109/3$&$-164$&180\\
4&$-1133/4$&907&$-715$&$-1489/4$&2226&$-3248$\\
5&$13856/5$&$-10936$&10632&$21771/5$&$-31468$&55684\\
6&$-58193/2$&134242&$-458096/3$&$-110719/2$&460910&$-2826242/3$\\
\bottomrule
\end{tabular}
\end{center}
The numbers of connected supports modulo translations, by sizes $1,\ldots,6$, are
\[
\begin{array}{c|rrrrrr}
\text{kings}&1&4&20&110&638&3832\\
\text{knights}&1&4&28&234&2162&20972.
\end{array}
\]
Inserting the table into~\eqref{eq:cformula} proves the explicit coefficients in~\eqref{eq:kings}--\eqref{eq:knights}. Formal division and logarithms give~\eqref{eq:ratio}--\eqref{eq:nlog}.

The first two uniform-density corrections, beyond~\eqref{eq:c1theta}, are
\begin{align*}
 c_1^{\rm K}(\theta)&=-\frac{\theta}{6}(49\theta^2-36\theta-27),\\
 c_2^{\rm K}(\theta)&=\frac{\theta}{72}
 (2401\theta^5-3528\theta^4-2628\theta^3
                    +4248\theta^2+891\theta-432),\\
 c_1^{\rm N}(\theta)&=-\frac{\theta}{6}(25\theta^2-72\theta-27),\\
 c_2^{\rm N}(\theta)&=\frac{\theta}{72}
 (625\theta^5-3600\theta^4+3924\theta^3
                    +7632\theta^2-405\theta-864).
\end{align*}

\subsection{Two neighboring board models}
For wazirs (orthogonal nearest-neighbor moves, A201511) and ferses (diagonal nearest-neighbor moves, A201861), the carrier is $e^{-5/2}n^{2n}/n!$. Their first four relative coefficients are
\[
\begin{array}{c|rrrr}
&c_1&c_2&c_3&c_4\\\hline
\text{wazirs}&7/3&32/9&-1003/810&68291/9720\\
\text{ferses}&13/3&128/9&8267/810&-388861/9720.
\end{array}
\]
These are instances of the same theorem, with the same all-fixed-order error guarantee.

\subsection{Finite computation and independent checks}
The accompanying \texttt{compute\_clusters.py} grows every connected support from the singleton and applies~\eqref{eq:weight} and~\eqref{eq:alphabetagamma}. Every finite connected graph has an ordering obtained by deleting leaves of a spanning tree, so every support occurs in this growth procedure. Translation normalization alone avoids erroneous symmetry division. Coefficients of graph logarithms are computed using the exact integer recurrence
\[
 jL_j=j i_j-\sum_{q=1}^{j-1}qL_q\,i_{j-q},
 \qquad I_H(z)=\sum_{j\ge0}i_jz^j.
\]
The file \texttt{derive\_expansion.py} implements~\eqref{eq:Hrec}--\eqref{eq:cformula} using rational symbolic algebra. No sequence fit is used.

The direct finite-board program \texttt{verify\_finite.py} enumerates independent subsets on $3\times3$, $4\times4$, and $5\times5$ boards for all four pieces. All twelve comparisons with the clipped-support formula~\eqref{eq:clipped} pass through cluster order six. A separate row-mask dynamic program reproduced both OEIS diagonals through $n=8$ and checked the cluster polynomials at all applicable board sizes through $13$ (111 exact comparisons). A separate Touchard-moment extraction verified the four displayed corrections for both main sequences. These finite checks validate the arithmetic implementation; Sections 3--5 supply the proof for all large $n$.

\section{Growth-index inversion}
Let $A_n$ be either main sequence. Stirling's formula gives, at every fixed order,
\begin{equation}\label{eq:logmodel}
 \log A_n=n(\log n+1)-\frac12\log(2\pi n)+a
          +\frac{d_1}{n}+\frac{d_2}{n^2}+\cdots,
\end{equation}
where $a=-9/2$ and
\[
 d_1=\frac94\quad\text{for kings},\qquad
 d_1=\frac{49}{4}\quad\text{for knights}.
\]
Replacing $n$ by a positive real variable in a fixed truncation defines a smooth asymptotic model; it does not define combinatorial counts on noninteger boards.

For $y\to\infty$, set
\begin{equation}\label{eq:inverse-def}
 L=\log y,\qquad x=\frac{L}{W(eL)},\qquad D=\log x+2,
 \qquad \delta_0=\frac{\tfrac12\log(2\pi x)-a}{D},
\end{equation}
where $W$ is the positive real Lambert function. For any logarithmic model retaining at least $d_1/n$, the inverse has
\begin{equation}\label{eq:inverse}
 n=x+\delta_0+
 \frac{\delta_0/2-\delta_0^2/2-d_1}{xD}
 +O\left(\frac1{x^2D}\right).
\end{equation}
Indeed $x(\log x+1)=L$. The first correction cancels the logarithmic and constant residual. Taylor expansion at $x+\delta_0$ leaves the term
\[
 \frac{\delta_0^2/2-\delta_0/2+d_1}{x}
\]
to first order in $x^{-1}$; division by the derivative $D+O(1/x)$ gives~\eqref{eq:inverse}. The omitted residual is $O(x^{-2})$. Recursive Taylor cancellation or Newton iteration gives every further fixed inverse order.

There is an essential qualification for integer thresholds. The leading equivalent proves eventual strict increase, since
\[
 \frac{A_{n+1}}{A_n}\sim e^2n.
\]
Let $q_M(y)$ solve the smooth logarithmic model through $d_M/n^M$. For all sufficiently large $y$, the threshold
\[
 N(y)=\min\{n\ge n_0:A_n\ge y\}
\]
is bracketed by the ceilings of
\begin{equation}\label{eq:round}
 q_M(y)\ \pm\ O\left(\frac{q_M(y)^{-M-1}}{\log q_M(y)}\right).
\end{equation}
This follows from the uniform pointwise logarithmic error and the positive derivative of the model. The two ceilings agree away from a shrinking window around integers. No finite-order expansion guarantees unconditional rounding when $y$ is arbitrarily close to a sequence value.

\section*{Scope of the literature check}
The OEIS definitions and cited book pages were inspected on 2 October 2026. Searches for these sequence IDs and higher-order chess expansions, and targeted checks of the current relevant ProveIt catalog, did not locate the specialized formulas displayed here. This is a bounded source check, not proof of absence from all literature. The general canonical cluster framework and the classical leading equivalents are explicitly credited. No convergent infinite asymptotic series or beyond-all-orders completion is claimed.

\begin{thebibliography}{9}
\bibitem{OEIS} The OEIS Foundation Inc.,
\href{https://oeis.org/A201513}{A201513},
\href{https://oeis.org/A201540}{A201540},
\href{https://oeis.org/A201511}{A201511}, and
\href{https://oeis.org/A201861}{A201861}. Accessed 2 October 2026.
\bibitem{Kote} V. Kot\v e\v sovec, \emph{Non-attacking Chess Pieces}, sixth edition, 2013, pp. 77, 293, 382, 423, 683, 685--686. The public file contains later table updates.
\href{https://www.kotesovec.cz/books/kotesovec_non_attacking_chess_pieces_2013_6ed.pdf}{Author's PDF}.
\bibitem{DJP} E. Davies, M. Jenssen, and W. Perkins,
\emph{A proof of the Upper Matching Conjecture for large graphs},
\href{https://arxiv.org/abs/2004.06695}{arXiv:2004.06695}, v2, 29 July 2021, especially Section 3.
\bibitem{PT} E. Pulvirenti and D. Tsagkarogiannis,
\emph{Cluster expansion in the canonical ensemble},
Communications in Mathematical Physics \textbf{316} (2012), 289--306.
\href{https://arxiv.org/abs/1105.1022}{arXiv:1105.1022}.
\end{thebibliography}
\end{document}
